\section{Scope, exact targets, and conventions}
\label{sec:scope}

The useful unit of progress in this programme is a proved identity or a
precisely delimited structural theorem.  Several nearby special-value
problems already have substantial proofs, and several plausible-looking
new directions are already present in earlier reports.  Accordingly,
the results below are attached to specific source questions before their
proofs are developed.  Bounds appear only inside the reconstruction and
convergence proofs; they are not independent objectives of the article.

The principal new conclusions are the complete order-eight theorem,
the effective fixed-order theorem, the all-moment quartic evaluation,
and the higher spectral factorization calculus.  These conclusions have
different logical content.  A nonzero rational determinant proves a
coefficient-rank statement; it does not prove independence of the real
constants later substituted into the coordinates.  A normally convergent
subtracted sum is an ordinary identity; the divergent raw sum preceding
that subtraction is not assigned an unspecified value.

\subsection{The concrete questions being continued}

\begin{longtable}{@{}>{\raggedright\arraybackslash}p{0.29\textwidth}>{\raggedright\arraybackslash}p{0.35\textwidth}>{\raggedright\arraybackslash}p{0.30\textwidth}@{}}
\toprule
Source target & Result here & Boundary retained\\
\midrule
\endhead
Reconstruction report, Q1: the omitted class
$b\equiv1\pmod{16}$ for $(8,2b)$ \cite{Reconstruction}
& Theorem~\ref{r8c:main} covers every $b\ge5$.
& The unrelated Gaussian period conjectures are unaffected.\\[5pt]
Reconstruction report, Q2: general unequal even orders
& Theorem~\ref{ert:main} gives a terminating threshold algorithm for
each fixed smaller order.
& The remaining finite set is not checked for every smaller order.\\[5pt]
Reconstruction report, Q5: a quartic formula, fixed space, and generator
& Theorems~\ref{qtm:main} and \ref{qtm:generator} give the complete
regular parity family and six formal coordinates.
& Fifth and higher powers require additional analysis.\\[5pt]
Quadratic-spectrum report, Q7: several factors beyond the first derivative
\cite{Quadratic}
& Theorems~\ref{msj:thm:second}, \ref{msj:thm:quadrature}, and
\ref{msj:thm:confluent} give the cross term and its full confluent series.
& A finite classical-polylogarithm reduction of the integral at arbitrary
parameters is not asserted.\\[5pt]
Quadratic-spectrum report, Q2: higher conjugate-linear defects
& Corollary~\ref{cor:linear-jets} gives every spectral order, with the
correct scale and all Stieltjes terms.
& The first-derivative parity criterion does not eliminate higher
nonlocal terms.\\[5pt]
Same-direction report, Q2: entire higher-depth connected interpolation
\cite{SameDirection}
& Section~\ref{sec:audit} identifies the existing all-depth interpolant
and geometric kernel exactly.
& This is an attribution and research-status correction, not a new
entire-family theorem.\\
\bottomrule
\end{longtable}

Two additional constructions go beyond those literal questions.
Theorem~\ref{thm:alljets} isolates the Hurwitz-pole contribution at every
derivative order for arbitrary polynomial weights and several quadratic
factors.  Theorem~\ref{thm:primitive} integrates a nonlinear second-jet
identity explicitly in the shift, using dilogarithms and a Gamma ratio.
Both have their normalization fixed in the statement.

\subsection{Source pin and attribution}

The repository was inspected at commit \sourcepin, including the
canonical manuscript status files and the relevant harmonic, Gaussian,
Hurwitz, and spectral branches.  All six archives currently in
\src{docs/incoming} at that pin were extracted and triaged.  The detailed
proof development uses the reconstruction, quadratic-spectrum, and
same-direction reports; the Beta-crossing, Gamma-continuation, and
shift-and-twist reports were also considered when selecting targets.
The exact source paths and SHA-256 digests are recorded in
\src{provenance.json}.  This is not an assertion that every theorem in
the four-volume manuscript has been audited.

Classical ingredients include Bernoulli--Hurwitz asymptotics, Euler's
Bernoulli formula, polygamma partial fractions, Gamma reflection, and
normalized Mellin continuation \cite{DLMF}.  Hoffman's quartic tail
sum \cite{Hoffman}, the Ihara--Nakamura--Yamamoto interpolant
\cite{INY}, and the first-derivative multiplicative-anomaly framework
\cite{EVZ,PolynomialAnomaly} are attributed where used.  Inherited
determinant reductions and cubic formulas are identified explicitly.
No claim of exhaustive historical priority is inferred from absence in
the inspected repository.

\subsection{Common conventions}

We use
\[
 H_n^{(s)}=\sum_{k=1}^n k^{-s},\qquad H_0^{(s)}=0,\qquad
 A_s(n)=\zeta(s,n+1),\qquad x_n=n+\tfrac12
\]
in the centered harmonic sections.  For an integer $s\ge2$,
$A_s(n)=\zeta(s)-H_n^{(s)}$.  The polygamma convention is
\[
 \psi^{(m)}(z)=(-1)^{m+1}m!\,\zeta(m+1,z),\qquad m\ge1.
\]
In particular $A_2(n)=\psi'(n+1)$ and
$A_{2k}(n)=\psi^{(2k-1)}(n+1)/(2k-1)!$.

Bernoulli polynomials are normalized by
\[
 \frac{te^{xt}}{e^t-1}=\sum_{m\ge0}B_m(x)\frac{t^m}{m!},
 \qquad B_m=B_m(0),\qquad B_1=-\tfrac12.
\]
The crucial centered zero is
\[
 \zeta(-2j,1/2)=0\qquad(j\ge0).
\]
It makes the specified polynomial subtractions have zero continued
sum.  A shifted analogue must restore its nonzero Hurwitz values;
this midpoint normalization is not silently transported to other shifts.

Superscripts on $\zeta$ and on spectral families denote derivatives
in their first, spectral variable unless stated otherwise.  We use
\[
 \zeta(1+u,a)=\frac1u+
       \sum_{m\ge0}\frac{(-1)^m}{m!}\gamma_m(a)u^m,
 \qquad \gamma_0(a)=-\psi(a).
\]
Thus the order derivative of $\zeta$ and the parameter derivative of a
Stieltjes constant are different operations.  All exchanges of spectral
differentiation and infinite sums below follow a locally uniform
subtraction argument, rather than formal differentiation at a pole.

The one-variable multiple polylogarithm uses the outer-index argument:
\[
 \Li_{s_1,\ldots,s_d}(z)=
 \sum_{n_1>\cdots>n_d\ge1}
       \frac{z^{n_1}}{n_1^{s_1}\cdots n_d^{s_d}}.
\]
At $z=1$ with admissible positive integer indices this is the
descending-index multiple zeta value.  In the finite generator formulas
we first take $z=e^{-t}$ with real $t>0$; the anchored integral formulas
specify the holomorphic continuation through the cancellations at $t=0$.

\subsection{Why the proofs use subtractions and leading determinants}

The harmonic and spectral problems share a finite algebraic step but have
different analytic mechanisms.  For reconstruction, the source matrix is
indefinite.  Its highest degree in a fixed number of Bernoulli variables
nevertheless factors into a positive discrete moment determinant.  The
factorial size of the Bernoulli variables makes that highest-degree part
dominate eventually.  The exact exceptional arithmetic remains visible
and is handled separately for order eight.

For quartic moments, summation by parts reduces the problem from the
fourth tail power to the inherited cubic coordinates.  A specific
boundary cancellation removes the remaining cubic Bernoulli convolution.
The surviving lower-power terms explain all exceptional low indices;
they cannot be discarded by parity alone.

For spectral derivatives, the common-scale difference vanishes to order
two in the spectral parameter before summation.  A resonant Hurwitz
pole can lower that order by one.  Its residue, regular Stieltjes
coefficients, and higher coefficient of the vanishing factor all matter.
The finite subtraction formulas retain those contributions while leaving
an ordinary summable remainder.  This is the source of the extra local
term in \eqref{one:x5second} and of the higher bilateral correction in
\eqref{msj:eq:third-contact}.
